\section{Diskussion}
\label{sec:Diskussion}

Die relative Messabweichung wird nach der Formel 
\begin{equation}
    \Delta x =   \frac{ \left| x_{exp} - x_{theo} \right|   }{x_{theo}}  
    \label{eq:abweichung}
\end{equation}
bestimmt. \\

Der Literaturwert der Lebensdauer eines Myons beträgt $\tau_{\text{Literatur}} = \qty{2.2}{\micro\second}$ \cite{Myon}.
Im Vergleich zum experimentell bestimmten Wert von $\tau_\T{exp} = \qty{2.48(0.35)}{\micro\second}$ ergibt sich eine relative Abweichung nach Gleichung \eqref{eq:abweichung} von $\Delta \tau = \qty{12.7(15.9)}{\%}$. \\
Der Theoriewert liegt im Bereich des Fehlers des experimentellen Wertes, wobei dazu gesagt werden muss, dass der Fehler des experimentell bestimmten Wertes ziemlich groß ist. Das liegt daran, dass nur wenige Stoppsignale gemessen wurden. Deshalb ist dieses Ergebnis auch ziemlich zufriedenstellend. Die Tendenz, dass die Lebenszeit der Myonen im Material etwas kleiner sein sollte als im Vakuum aufgrund des Myoneneinfangs lässt sich hier nicht beobachten, da aufgrund der geringen Anzahl an Stoppsignalen eine zu hohe Unsicherheit herrscht. \\
Dass es nur so wenige Stoppsignale gibt liegt daran, dass von vornerein die Anzahl der Signale die durch die Koinzidenzschaltung durchgegangen sind schon ziemlich niedrig waren. Die dort gemessene Impulsrate betrug auch nach dem Einstellen der Verzögerung eher $\qty{15}{\hertz}$ als die geforderten $\qty{20}{\hertz}$. Des Weiteren wurde nach einer Messzeit von zwei Tagen bemerkt, dass es viel zu wenige Stoppsignale gibt, weshalb die Spannung des Photomultpliers hochgedreht wurde. Über den Messzeitraum des dritten Tages kamen dann ca. doppelt so viele Stoppsignale an, insgesamt aber trotzdem eher wenige Stoppsignale mit ca. $\qty{400}{d^{-1}}$.
Außerdem fiel beim Ablesen der Start- und Endimpulse auf, dass zwar der Zähler für die Stoppimpulse weiterlief, der für die Startimpulse aber nicht, dieser blieb bei circa $N_{\text{Start}}=980000$ stehen. 
Der ursprüngliche Grund dafür, dass erst so wenige Signale die Koinzidenzschaltung erreichten, könnte sein, dass der Versuch ziemlich alt ist und deswegen das Szinitllatormaterial teilweise nicht mehr so effizient funktionieren könnte.
Auch der Diskriminator könnte ein entscheidender Faktor dafür sein, dass es zu wenige Stoppsignale gibt, denn dieser wurde eigentlich so eingestellt, dass sein Threshhold sehr niedrig ist und nahezu jedes Signal durchlassen müsste. Hierbei wurde aber über die Schraube zur Einstellung des Thresholds über das Minimum hinaus gedreht, wodurch der Diskriminator unvorhersehbare Dinge tun könnte.

\clearpage